\documentclass[12pt]{article}

\usepackage[utf8]{inputenc}
\usepackage[T1]{fontenc}
\usepackage[brazil]{babel}

\usepackage{geometry}
\geometry{margin=2.5cm}
\usepackage{setspace}
\onehalfspacing

\usepackage{amsmath, amssymb, amsthm}
\usepackage{mathtools}

\numberwithin{equation}{section}

\theoremstyle{plain}
\newtheorem{theorem}{Teorema}[section]
\newtheorem{proposition}[theorem]{Proposição}
\newtheorem{lemma}[theorem]{Lema}
\newtheorem{corollary}[theorem]{Corolário}

\theoremstyle{definition}
\newtheorem{definition}[theorem]{Definição}
\newtheorem{example}[theorem]{Exemplo}
\newtheorem{exercise}{Exercício}[section]

\theoremstyle{remark}
\newtheorem*{remark}{Observação}

\usepackage{hyperref}
\usepackage{csquotes}

\usepackage[
    backend=biber,
    style=authoryear,
    maxcitenames=2
]{biblatex}

\addbibresource{/mnt/c/Users/nicho/Documents/nicholasvoltani.github.io/bibliography.bib}

\newcommand{\R}{\mathbb{R}}
\newcommand{\pref}{\succeq}
\newcommand{\strictpref}{\succ}
\newcommand{\indiff}{\sim}

% ------------------------------------------------
% Title (fill manually)
% ------------------------------------------------
\title{}
\author{}
\date{}

\begin{document}

\maketitle

\hypertarget{cobb-douglas-para-l-bens}{%
\section{\texorpdfstring{Cobb-Douglas para \(L\)
bens}{Cobb-Douglas para L bens}}\label{cobb-douglas-para-l-bens}}

A Função de Cobb-Douglas geral tem a forma \[
U(x_{1},\dots,x_{L}) = \prod \limits_{k=1}^{L} x_{k}^{\alpha_{k}}
\] Suponhamos que ela já seja normalizada, de tal forma que
\(\sum_{k} \alpha_{k} = 1\) (com \(\alpha_{k} > 0\)).

Tenha-se uma Restrição Orçamentária da forma
\(\braket{ p | x } \leq w\).

Busque-se uma \textit{solução interior}. Então o lagrangiano associado é
\[
\mathcal{L}(x, p, w) = \prod \limits_{k=1}^{L} x_{k}^{\alpha_{k}} + \lambda \left(\braket{ p | x } - w\right)
\]

As condições de primeira ordem dão \[
\begin{cases}
\forall k: \frac{\alpha_{k}}{x_{k}} \prod \limits_{k=1}^{L} x_{k}^{\alpha_{k}}  + \lambda p_{k} &= 0\\
\sum \limits_{k} p_{k} x_{k} &= w
\end{cases}
\]

Das primeiras equações, tem-se que, para quaisquer \(i, j\), vale \[
\frac{\alpha_{i}}{p_{i} x_{i}} = \frac{\alpha_{j}}{p_{j} x_{j}}
\]

Da restrição, obtém-se que, para qualquer \(i\) \[
\begin{align}
w = \sum\limits_{k} \frac{p_{k} x_{k}}{\alpha_{k}} \alpha_{k} &= \frac{p_{i} x_{i}}{\alpha_{i}} \sum \limits_{k} \alpha_{k}  \\
&= \frac{p_{i} x_{i}}{\alpha_{i}}
\end{align}
\] Portanto, para todo \(i\) vale \[
x_{i} = \alpha_{i} \frac{w}{p_{i}}
\] Equivalentemente, vale o resultado esperado para uma utilidade
Cobb-Douglas: \[
\alpha_{i} = \frac{p_{i} x_{i}}{w}
\] Ou seja, os argumentos \(\alpha_{i}\) dão as \textit{proporções} de
cada bem \(i\) no dispêndio da renda \(w\){[}\^{}1{]}.

\hypertarget{matriz-de-slutsky-para-utilidades-e-demandas-conhecidas}{%
\section{Matriz de Slutsky para utilidades (e demandas)
conhecidas}\label{matriz-de-slutsky-para-utilidades-e-demandas-conhecidas}}

\hypertarget{cobb-douglas}{%
\subsection{Cobb-Douglas}\label{cobb-douglas}}

\hypertarget{ces}{%
\subsection{CES (?)}\label{ces}}

\printbibliography

\end{document}
